\subsection{Intersectional Subgroup Performance}
\label{sec:intersectional_subgroup_performance}
\label{sec:subgroup_performance}

The granular OvR classification performance profiles of the best tri-domain models (LR and SVM) are evaluated across six demographic subgroups ($N = 360$ test images per subgroup), as presented in Table~\ref{tab:15}. Across both models, the White Males subgroup records the highest performance, with a Recall of 94.44\% (F1-Score of 94.84\% in LR and 94.18\% in SVM) and the lowest False Negative Rate (FNR) of 5.56\%. Conversely, the Black Females subgroup records a Recall of 84.72\% (F1-Score of 84.96\% in LR and 86.16\% in SVM) with an FNR of 15.28\%. The Asian Females subgroup achieves a Recall of 86.67\% in LR and 86.11\% in SVM. All subgroup cohorts maintain an F1-Score above 84.90\%, confirming the reliability of tri-domain feature fusion in separating intersectional demographic boundaries.

\begin{table*}[htbp]
\caption{Intersectional Subgroup Performance and Algorithmic Fairness Profile.}
\label{tab:15}
\label{tab:XV}
\centering
\footnotesize
\setlength{\tabcolsep}{6pt}
\resizebox{\textwidth}{!}{%
\begin{tabular}{llccccc}
\toprule
\textbf{Classifier} & \textbf{Subgroup} & \textbf{Recall (TPR)} & \textbf{Precision} & \textbf{F1-Score} & \textbf{OvR Accuracy} & \textbf{FNR} \\
\midrule
\multirow{6}{*}{Logistic Regression (Tri-Domain)} 
 & White Males   & 94.44\% & 95.24\% & 94.84\% & 98.29\% & 5.56\% \\
 & White Females & 90.56\% & 91.32\% & 90.93\% & 96.99\% & 9.44\% \\
 & Black Males   & 87.78\% & 91.59\% & 89.65\% & 96.62\% & 12.22\% \\
 & Asian Males   & 87.78\% & 82.72\% & 85.18\% & 94.91\% & 12.22\% \\
 & Asian Females & 86.67\% & 86.43\% & 86.55\% & 95.51\% & 13.33\% \\
 & Black Females & 84.72\% & 85.20\% & 84.96\% & 95.00\% & 15.28\% \\
\midrule
\multirow{6}{*}{Support Vector Machine (Tri-Domain)} 
 & White Males   & 94.44\% & 93.92\% & 94.18\% & 98.06\% & 5.56\% \\
 & White Females & 91.11\% & 90.11\% & 90.61\% & 96.85\% & 8.89\% \\
 & Asian Males   & 88.06\% & 82.34\% & 85.10\% & 94.86\% & 11.94\% \\
 & Black Males   & 87.22\% & 92.35\% & 89.71\% & 96.67\% & 12.78\% \\
 & Asian Females & 86.11\% & 85.87\% & 85.99\% & 95.32\% & 13.89\% \\
 & Black Females & 84.72\% & 87.64\% & 86.16\% & 95.46\% & 15.28\% \\
\bottomrule
\end{tabular}%
}
\end{table*}

The evaluation of algorithmic fairness disparities across subgroups is quantified using the EOD metric based on the difference between the maximum and minimum TPR ($\text{EOD} = \Delta \text{TPR} = \max \text{TPR} - \min \text{TPR} = 94.44\% - 84.72\% = \mathbf{9.72\text{ pp}}$), which is identical for both the LR and SVM models. This disparity reflects variations in the empirical error frequency distribution across subgroups on the held-out test data, where Black Females record the highest FNR (15.28\%) and White Males the lowest (5.56\%). As a commitment to Artificial Intelligence (AI) ethics and scope delimitation (Responsible AI), this disparity is explicitly not claimed as a guarantee of achieving universal demographic fairness (no overclaim of complete demographic fairness), but rather as an empirical subgroup performance transparency audit.
